\documentclass[conference]{IEEEtran}
\usepackage{cite}
\usepackage{amsmath,amssymb,amsfonts}
\usepackage{graphicx}
\usepackage{textcomp}
\usepackage{xcolor}
\usepackage{booktabs}
\usepackage{multirow}
\usepackage{url}
\usepackage{hyperref}
\usepackage{array}

\begin{document}

\title{S2S: Autonomous Multi-Agent Specification-to-GDSII\\Chip Design with Deterministic Verification}

\author{
\IEEEauthorblockN{Goutham Rajam}
\IEEEauthorblockA{Independent Researcher\\
Repository: \url{https://github.com/GouthamRajam/Gouthamlabs}}
}

\maketitle

\begin{abstract}
We present S2S (Spec-to-Silicon), a multi-agent framework that autonomously transforms a chip specification into tape-out-ready GDSII using open-source EDA tools and SkyWater SKY130. Thirty-seven AI agents spanning seven disciplines are coordinated by a Trust Layer enforcing a strict invariant: \emph{real tools decide pass/fail---the AI never does}. No stage advances unless a deterministic tool (simulator, DRC, LVS, STA) runs and passes. We evaluate on seven designs from an 8-bit counter to a RISC-V RV32I core. Four of five Phase~0 designs reach DRC-clean GDSII; a PicoRV32 core passes 39/39 compliance tests and achieves signoff across 9~PVT corners at 20~MHz. The autonomous pipeline produces a complete UART---specification through routed layout---at \$0.39 and 3.5~minutes. We report all results honestly, including failures, a provenance audit that corrected false attribution, and every known limitation.
\end{abstract}

\begin{IEEEkeywords}
chip design automation, multi-agent systems, LLM, EDA, RISC-V, hardware verification, open-source EDA
\end{IEEEkeywords}

% ===========================================================================
\section{Introduction}
\label{sec:intro}

Designing a digital chip from specification to GDSII requires specialists working years, with verification consuming the majority of effort. First-silicon success rates remain approximately 14\%~\cite{wilson2024}, and each respin costs millions. Meanwhile, demand for custom silicon grows across AI, automotive, and IoT while the talent shortage worsens~\cite{semiengineering2023}.

Recent work applies LLMs to individual stages: RTL generation~\cite{thakur2024verigen, liu2023chipnemo}, verification~\cite{orenesvera2023using}, and design exploration. However, no prior system demonstrates complete, autonomous specification-to-GDSII with deterministic verification at every stage.

We present \textbf{S2S} with three contributions:
\begin{enumerate}
    \item \textbf{Trust Layer}: A deterministic verification invariant preventing hallucination propagation---no stage advances unless a real EDA tool runs and passes.
    \item \textbf{Full-flow autonomy}: 37 agents covering spec through signoff, producing GDSII from text, with verification agents structurally independent from RTL agents.
    \item \textbf{Reproducible evidence}: Trajectory logs, tool reports, and signoff evidence for seven designs including a RISC-V core, using exclusively open-source tools on SKY130.
\end{enumerate}

% ===========================================================================
\section{Related Work}
\label{sec:related}

\textbf{LLM-assisted hardware.} VeriGen~\cite{thakur2024verigen} fine-tunes models on Verilog. ChipNemo~\cite{liu2023chipnemo} adapts foundation models for chip design tasks. AutoChip~\cite{thakur2023autochip} applies iterative LLM feedback with simulation correction. These focus on individual stages and do not address physical implementation or signoff.

\textbf{Full-flow systems.} ASIC-Agent~\cite{asicagent2025} demonstrates multi-agent RTL-to-Caravel integration with a benchmark suite. NL2GDS~\cite{nl2gds2026} pairs an LLM front end with OpenLane for natural-language-to-GDSII. AutoEDA~\cite{autoeda2025} uses microservice-based agents for RTL-to-GDSII automation. A comprehensive survey appeared in late 2025~\cite{agenticseda2025}.

S2S differs in three ways: (1)~the Trust Layer's 5-condition safety net is a formal invariant, not iterative debugging; (2)~verification agents are structurally independent from RTL agents (filesystem-enforced); (3)~we report full signoff evidence with honest failure analysis, including designs that did not pass.

\textbf{Open-source EDA.} The OpenLane~\cite{openlane} flow (Yosys~\cite{yosys}, OpenROAD~\cite{openroad}, Magic~\cite{magic}, KLayout~\cite{klayout}, Netgen~\cite{netgen}) enables complete RTL-to-GDSII on SKY130~\cite{sky130}. S2S uses these as its deterministic backbone.

% ===========================================================================
\section{Framework Architecture}
\label{sec:architecture}

\subsection{Design Flow and Agents}

S2S organizes chip design into 10 stages in four zones (Fig.~\ref{fig:flow}): Z1 (specification, architecture), Z2A (RTL implementation), Z2B/2C (independent test generation and verification closure), Z3 (synthesis through P\&R), and Z4 (signoff and GDSII). Zones 2A and 2B execute in parallel from the same specification; the \emph{independence constraint} prevents shared-mistake propagation.

\begin{figure}[t]
\centering
\fbox{\parbox{0.9\columnwidth}{\small
\textbf{Z1:} Spec $\rightarrow$ Architecture $\rightarrow$ \\
\textbf{Z2A:} RTL \quad $\|$ \quad \textbf{Z2B:} Tests (independent) \\
$\downarrow$ \\
\textbf{Z2C:} Verification closure \\
$\downarrow$ \\
\textbf{Z3:} Synth $\rightarrow$ Floorplan $\rightarrow$ P\&R \\
$\downarrow$ \\
\textbf{Z4:} Signoff (DRC/LVS/STA/Antenna) $\rightarrow$ GDSII
}}
\caption{S2S design flow. Zones 2A and 2B execute in parallel; verification agents cannot read RTL (filesystem-enforced).}
\label{fig:flow}
\end{figure}

The 37 agents span seven disciplines: specification (4), architecture (3), RTL (4), verification (8), physical design (5), cross-cutting (8), and knowledge (5). Each has defined scope, tool permissions, and budget limits. Key constraints include: verification agents cannot read \texttt{rtl/**} (filesystem-enforced), N-version reference models flag disagreement, and an anti-cheating guard detects test weakening via AST diffs.

\subsection{The Trust Layer}

The Trust Layer enforces: \emph{a stage advances only when a deterministic tool runs and passes}. Five conditions must all hold:
\begin{enumerate}
    \item Tool exit code is 0
    \item Zero error markers in output
    \item No stage-specific rule violated (no latches, no logic loops, no negative slack)
    \item No known failure strings
    \item Positive success marker present
\end{enumerate}

This catches tools that report errors but return exit code 0---common EDA behavior. Three consecutive failures trigger human escalation; budget exhaustion pauses the task cleanly.

\subsection{Tool Integration}

Ten adapters wrap open-source tools, each in a Docker container (\texttt{--rm --network none}): Verilator (lint), Yosys (synthesis), OpenSTA (timing), Icarus Verilog and cocotb (simulation), SymbiYosys (formal), eqy (equivalence), OpenROAD/OpenLane (P\&R), Magic and KLayout (DRC), and Netgen (LVS).

% ===========================================================================
\section{Experimental Setup}
\label{sec:setup}

We evaluate on seven designs (Table~\ref{tab:designs}) targeting SKY130 \texttt{sky130\_fd\_sc\_hd}. Timing signoff uses 9~PVT corners with OCV derates at $\pm$5\% and $\pm$10\%. Agents use Claude Sonnet with per-agent budgets (300K--400K tokens, 16--40 tool calls). The orchestrator operates as deterministic rules.

\begin{table}[t]
\centering
\caption{Design Targets}
\label{tab:designs}
\begin{tabular}{@{}llll@{}}
\toprule
\textbf{Design} & \textbf{Type} & \textbf{Phase} & \textbf{Freq.} \\
\midrule
counter8 & 8-bit counter & 0 & 50 MHz \\
alu8 & 8-bit ALU & 0 & 50 MHz \\
uart & UART transceiver & 0 & 50 MHz \\
sync\_fifo & Synchronous FIFO & 0 & 50 MHz \\
spi\_master & SPI master & 0 & 50 MHz \\
picorv32 & RV32I CPU & 1 & 20 MHz \\
\bottomrule
\end{tabular}
\end{table}

% ===========================================================================
\section{Results}
\label{sec:results}

\subsection{Phase 0: Five Digital Blocks}

\begin{table*}[t]
\centering
\caption{Phase 0 Signoff Results on SKY130}
\label{tab:phase0}
\begin{tabular}{@{}lrrrrrrrl@{}}
\toprule
\textbf{Design} & \textbf{Cells} & \textbf{Area ($\mu$m$^2$)} & \textbf{DRC} & \textbf{LVS} & \textbf{Ant.} & \textbf{STA} & \textbf{Mut.} & \textbf{Status} \\
\midrule
counter8 & 71 & 3,119 & 0 & Match & 0 & MET & 100\% & PASS \\
alu8 & 156 & 5,478 & 0 & Match & 0 & MET & --- & PASS \\
uart & 208 & 7,834 & 0 & Match & 0 & MET & 95\% & PASS \\
spi\_master & 312 & 11,267 & 0 & Match & 0 & MET* & 85\% & COND. \\
sync\_fifo & 189 & 6,912 & 0 & Match & 2 & MET & 100\% & FAIL \\
\bottomrule
\multicolumn{9}{@{}l@{}}{\footnotesize *CTS fanout relaxed 10$\to$20. sync\_fifo: 2 antenna violations (13--17\% over limit).} \\
\end{tabular}
\end{table*}

Four designs achieve clean GDSII. The uart required 4~P\&R iterations (24~tool runs); the Trust Layer caught DRC violations after each attempt, preventing bad layouts from advancing. Across all designs, the safety-net parser caught 71~false-stall alarms from ambiguous tool output.

\subsection{Failure Analysis}

\textbf{sync\_fifo (FAIL):} Default synthesis produced $-1.545$~ns slack; switching to DELAY~4 recovered to $+0.182$~ns at 19\% more area. After routing, 2~nets had antenna ratios 13--17\% above limits. Root cause: OpenROAD's post-DRT antenna repair did not converge on the pinned version (OpenLane~2.3.10), and pre-DRT diode insertion was insufficient because routed topology differed from global-route estimates. The Trust Layer blocked GDSII generation; the anti-weakening guard prevented any threshold relaxation.

\textbf{spi\_master (CONDITIONAL):} Passed all checks except clock-net fanout (CTS buffers drove $>$10 loads, relaxed to 20---a legitimate trade-off, as CTS buffer datasheets specify 20--30 max fanout). Formal verification achieved only 85\% mutation kill rate; 15\% survivors concentrated in SPI mode-switching logic where agents lacked protocol context.

\subsection{Phase 1: RISC-V RV32I Core}

\begin{table}[t]
\centering
\caption{PicoRV32 Results (SKY130, 20 MHz)}
\label{tab:phase1}
\begin{tabular}{@{}ll@{}}
\toprule
\textbf{Metric} & \textbf{Value} \\
\midrule
RV32I compliance & 39/39 PASS (Sail v0.15) \\
Cells / Die area & 8,374 / 297,307 $\mu$m$^2$ \\
DRC / LVS / Antenna & 0 / Match / 0 \\
STA (9 corners) & WNS = 0 (setup \& hold) \\
Antenna diodes & 11,631 (44.3\% of instances) \\
Power (VCD / stat.) & 2.606 / 3.387 mW \\
IR drop & 0.01\% of VDD \\
ATPG coverage & 83.1\% (random only) \\
\bottomrule
\end{tabular}
\end{table}

The PicoRV32 passes all checks, but the conservative antenna heuristic (diode on every net $>$50\% of limit, pre-DRT) inserted 11,631 diodes---44.3\% of instances. A post-DRT approach using actual wire lengths would reduce this significantly. The 83.1\% ATPG uses random patterns only; commercial stuck-at tools would achieve higher coverage.

\subsection{Autonomous Pipeline UART}

The pipeline-generated UART (run \texttt{01M426JH3Y...}) demonstrates fully autonomous operation: 12~model calls, zero operator influence, by agents Ada, Ada-C, Ada-I (RTL) and Lumen (lint).

\begin{table}[t]
\centering
\caption{Pipeline UART: Autonomous Spec-to-Layout}
\label{tab:pipeline_uart}
\begin{tabular}{@{}ll@{}}
\toprule
\textbf{Stage} & \textbf{Result} \\
\midrule
Verification & 72/72 loopback + reference model \\
Synthesis & 474 cells, 4,079 $\mu$m$^2$ \\
P\&R & 95,375 wires, 2 antenna repaired \\
STA post-P\&R & Setup +9.11~ns, hold +0.52~ns \\
Cost / Time & \$0.39 / 3.5 min \\
\bottomrule
\end{tabular}
\end{table}

Success factors: well-constrained spec (standard UART protocol), independent verification catching a bit-ordering bug before synthesis, and small design closing timing on the first P\&R iteration. Two antenna violations found during global routing were successfully repaired with diode insertion before detailed routing---the same approach that failed on sync\_fifo worked here because the routing topology was simpler (474 cells vs.\ 189 with deeper pointer logic paths).

\subsection{Verification Organization}

The verification methodology achieves mutation testing kill rates of 100\% (counter8, sync\_fifo), 95\% (uart), and 85\% (spi\_master) using 6~bug types and 20~mutants per design, type-balanced. Formal properties (7--14 per design) are proven via k-induction in SymbiYosys with vacuity and tautology checks. Pre/post-synthesis equivalence is verified via eqy (BMC, depth 10). A prove-then-exclude policy uses eqy to screen out provably equivalent mutants before counting, correcting uart's kill rate from 70\% to 95\% after identifying 8~equivalent mutants.

\subsection{Pipeline Statistics}

Across 25~pipeline runs: 5,630 trajectory records, 337 model calls, 7 designs. Twelve runs were fully autonomous; thirteen required intervention. Failure modes: token budget exhaustion (5), formal harness bugs (3), container timeouts (3), equivalent mutant miscounting (2). Each was logged, root-caused, and fed into the knowledge base.

% ===========================================================================
\section{Provenance and Limitations}
\label{sec:provenance}

\textbf{Provenance audit.} A provenance audit revealed that committed RTL for s2s\_rv32i and s2s\_uart was hand-authored in interactive sessions, not agent-generated. False labels were corrected. The pipeline does function---proven by the UART trajectory with SHA-256 hashes---but the two processor-scale designs were human-directed.

\textbf{Known limitations:} (1)~Bounded equivalence only (BMC depth 10, not full k-induction). (2)~Post-DRT antenna repair does not converge; pre-DRT heuristic is the workaround. (3)~Agent persona descriptions are prompting devices, not real expertise. (4)~Open-shuttle signoff bar, not commercial production. (5)~83.1\% ATPG from random patterns only. (6)~44.3\% diode overhead from conservative heuristic. (7)~No commercial-node validation.

% ===========================================================================
\section{Discussion}
\label{sec:discussion}

The Trust Layer addresses the fundamental challenge of LLM hallucination in safety-critical engineering: you don't trust the AI---you trust the tools. Every stage has a deterministic checker the AI cannot bypass. This pattern generalizes to any domain with deterministic verification (compilers, structural engineering, drug interaction checking).

The independence requirement---verification agents cannot read RTL, enforced by filesystem permissions---applies N-version programming~\cite{avizienis1985nversion} to AI-generated hardware. When RTL and tests derive from the same spec by different agents, they are unlikely to share mistakes.

At \$0.39 per design and 3.5~minutes, the pipeline enables design-space exploration at a scale impractical with human-only teams. A designer could specify dozens of architectural variants and let the pipeline evaluate each through synthesis and timing analysis. The 5,630 trajectory records with 337 model calls---each including agent prompt, response, tool calls, and outcomes---constitute a dataset for studying AI-driven chip design, with 13 runs classified as ``training-grade'' (complete provenance) and 12 as ``usable.''

% ===========================================================================
\section{Conclusion}
\label{sec:conclusion}

S2S demonstrates that LLM-powered agents can produce functionally correct, physically implementable digital designs when coupled with a Trust Layer delegating all pass/fail decisions to deterministic tools. Four of five Phase~0 designs reach clean GDSII; a RISC-V core passes 39/39 compliance tests across 9~PVT corners; the autonomous pipeline produces a UART at \$0.39 and 3.5~minutes. All results, logs, and signoff evidence are openly available.

Future work: mixed-signal extension, commercial-node validation, improved antenna repair, unbounded equivalence proofs, and first-silicon fabrication on an open shuttle.

\section*{Data Availability}

Codebase, agents, trajectory logs, and signoff evidence: \url{https://github.com/GouthamRajam/Gouthamlabs}.

\section*{Acknowledgments}

This work used exclusively open-source EDA tools and SKY130. The author thanks the OpenROAD, Yosys, OpenLane, Magic, KLayout, and Netgen communities. AI assistance provided by Anthropic's Claude.

% ===========================================================================
\begin{thebibliography}{20}

\bibitem{wilson2024}
H.~Foster, ``2024 Wilson Research Group Functional Verification Study,'' Siemens EDA, 2024.

\bibitem{semiengineering2023}
M.~LaPedus, ``Chip engineering talent shortage worsens,'' Semiconductor Engineering, 2023.

\bibitem{thakur2024verigen}
S.~Thakur \emph{et al.}, ``VeriGen: A large language model for Verilog code generation,'' \emph{ACM TODAES}, vol.~29, no.~3, 2024.

\bibitem{liu2023chipnemo}
M.~Liu \emph{et al.}, ``ChipNemo: Domain-adapted LLMs for chip design,'' in \emph{Proc. ISCA}, 2024.

\bibitem{orenesvera2023using}
B.~Orenes-Vera \emph{et al.}, ``Using LLMs for hardware verification,'' in \emph{Proc. MICRO Workshop}, 2023.

\bibitem{thakur2023autochip}
S.~Thakur \emph{et al.}, ``AutoChip: Automating HDL generation using LLM feedback,'' in \emph{Proc. DAC Workshop}, 2023.

\bibitem{asicagent2025}
M.~Allam \emph{et al.}, ``ASIC-Agent: Autonomous multi-agent ASIC design with benchmark,'' \emph{arXiv:2508.15940}, 2025.

\bibitem{nl2gds2026}
T.~Eland \emph{et al.}, ``NL2GDS: LLM-aided interface for open source chip design,'' \emph{arXiv:2603.05489}, 2026.

\bibitem{autoeda2025}
``AutoEDA: Enabling EDA flow automation through microservice-based LLM agents,'' \emph{arXiv:2508.01012}, 2025.

\bibitem{agenticseda2025}
``The dawn of agentic EDA: A survey of autonomous digital chip design,'' \emph{arXiv:2512.23189}, 2025.

\bibitem{avizienis1985nversion}
A.~Avizienis, ``The N-version approach to fault-tolerant software,'' \emph{IEEE TSE}, vol.~SE-11, no.~12, 1985.

\bibitem{openlane}
M.~Shalan and T.~Edwards, ``Building OpenLane: A 130nm tapeout-proven flow,'' in \emph{Proc. DAC Workshop}, 2020.

\bibitem{yosys}
C.~Wolf, ``Yosys Open SYnthesis Suite.'' [Online]. Available: \url{https://yosyshq.net/yosys/}

\bibitem{openroad}
A.~Kahng \emph{et al.}, ``OpenROAD: Open-source RTL-to-GDS flow,'' in \emph{Proc. DAC}, 2021.

\bibitem{magic}
R.~T.~Edwards, ``Magic VLSI layout tool.'' [Online]. Available: \url{http://opencircuitdesign.com/magic/}

\bibitem{klayout}
M.~Koefferlein, ``KLayout.'' [Online]. Available: \url{https://www.klayout.de/}

\bibitem{netgen}
R.~T.~Edwards, ``Netgen: LVS for VLSI.'' [Online]. Available: \url{http://opencircuitdesign.com/netgen/}

\bibitem{sky130}
Google and SkyWater, ``SkyWater SKY130 open-source PDK,'' 2020.

\bibitem{picorv32}
C.~Wolf, ``PicoRV32 -- A size-optimized RISC-V CPU.'' [Online]. Available: \url{https://github.com/YosysHQ/picorv32}

\bibitem{anthropic2024claude}
Anthropic, ``Claude,'' 2024. [Online]. Available: \url{https://www.anthropic.com/claude}

\end{thebibliography}

\end{document}
